%----------------------------------------------------------------------------------------
% STELSELS EN DE METHODE VAN GAUSS-JORDAN. Alle opgaven zijn eigen werk.
%
% Opbouw: Start, leerstof (theorie, werkwijze of uitleg, voorbeeld, korte oefening), en op een
% nieuwe pagina de Opdrachten (door elkaar, stijgende moeilijkheid), Koppel en Test jezelf.
% Labels van dit hoofdstuk beginnen met h2-.
%----------------------------------------------------------------------------------------

\mychapter{2}{Stelsels en de methode van Gauss-Jordan}

\lkr{Leerplan B+S'': LPD 37 (stelsels van eerstegraadsvergelijkingen oplossen met de methode van Gauss-Jordan), met de wenken: minstens één stelsel met drie onbekenden, aandacht voor bepaalde, onbepaalde en strijdige stelsels, en stelsels met één parameter. Rekenwerk met de hand blijft beperkt; bij moeilijkere stelsels mag het rekentoestel (op de TI-84: \texttt{rref(} in het menu MATRIX, MATH). De regel van Cramer komt in hoofdstuk 3, bij de determinant.}

In de 2de graad loste je stelsels van twee vergelijkingen met twee onbekenden op, door een onbekende te elimineren of te substitueren. Met drie of meer onbekenden wordt dat snel onoverzichtelijk. In dit hoofdstuk leer je een vaste methode die altijd werkt: je schrijft alleen de getallen van het stelsel in een matrix, en je vereenvoudigt die rij per rij tot je de oplossing kunt aflezen. Zo zie je ook meteen of een stelsel één oplossing heeft, oneindig veel, of geen enkele.

\begin{oefeningen}[Start]
\oef{h2-start-1}

\oef{h2-verken}
\end{oefeningen}

\lkr{Laat oefening \ref{oef:h2-verken} maken vóór de uitleg. In a lossen leerlingen een stelsel in trapvorm op van onder naar boven. In b moeten ze zelf zo'n trapvorm maken door $x$ te elimineren uit de tweede en de derde vergelijking. In c benoemen ze de toegelaten bewerkingen: dat zijn de rijbewerkingen van sectie 2.}

%========================================================================================
\section{Stelsels en matrices}

Een stelsel van eerstegraadsvergelijkingen heet ook een \textit{lineair} stelsel: de onbekenden komen alleen voor in de eerste macht, vermenigvuldigd met een getal. Alles wat je nodig hebt, zijn die getallen. Met het product van matrices uit hoofdstuk 1 schrijf je het hele stelsel als één vergelijking.

\begin{theorie}
Een stelsel van $m$ eerstegraadsvergelijkingen met $n$ onbekenden schrijf je als \quad $\formule{AX = B}$

\vspace{1mm}
Hierin is $A$ de $m \times n$-\textbf{\textcolor{accent}{coëfficiëntenmatrix}}, $X$ de kolom met de onbekenden en $B$ de kolom met de constanten. De \textbf{\textcolor{accent}{uitgebreide matrix}} van het stelsel is $\formule{[A \,|\, B]}$: de matrix $A$ met $B$ als extra kolom.
\end{theorie}

\voorbeeld{}{Schrijf het stelsel $\begin{cases} 3x - y + 2z = 5 \\ x + 4z = -1 \\ 2y - z = 0 \end{cases}$ als matrixvergelijking, en geef de uitgebreide matrix.\label{vb:h2-matrixvorm}}{%
\uitleg{Elke vergelijking wordt een rij, elke onbekende een kolom. Ontbreekt een onbekende in een vergelijking, dan is haar coëfficiënt 0.}
\begin{stappenbreed}
\begin{bmatrix} 3 & -1 & 2 \\ 1 & 0 & 4 \\ 0 & 2 & -1 \end{bmatrix} \begin{bmatrix} x \\ y \\ z \end{bmatrix} = \begin{bmatrix} 5 \\ -1 \\ 0 \end{bmatrix} & Rij maal kolom geeft telkens het linkerlid: $3x - y + 2z$, enzovoort. In de tweede vergelijking ontbreekt $y$: een 0. \\[6ex]
\um{ccc|c}{3 & -1 & 2 & 5 \\ 1 & 0 & 4 & -1 \\ 0 & 2 & -1 & 0} & De uitgebreide matrix. De streep staat waar het gelijkheidsteken stond.
\end{stappenbreed}}

\begin{oefeningen}
\oef{h2-matrixvorm}
\end{oefeningen}

%========================================================================================
\section{De methode van Gauss-Jordan}

\subsection{Rijbewerkingen}

In de Verken heb je vergelijkingen van elkaar afgetrokken om een onbekende te laten verdwijnen. Op de uitgebreide matrix zijn dat bewerkingen met rijen. Drie soorten zijn toegelaten, omdat ze de oplossingen van het stelsel niet veranderen.

\begin{theorie}
De \textbf{\textcolor{accent}{elementaire rijbewerkingen}} zijn:

\vspace{1mm}
{\centering\arrayrulecolor{accent}\renewcommand{\arraystretch}{1.5}%
\begin{tabular}{|>{\centering\arraybackslash}m{6.6cm}|>{\centering\arraybackslash}m{3.6cm}|}
\hline
twee rijen van plaats verwisselen & $\formule{R_i \leftrightarrow R_j}$ \\ \hline
een rij vermenigvuldigen met een getal $r \neq 0$ & $\formule{R_i \to rR_i}$ \\ \hline
bij een rij een veelvoud van een andere rij optellen & $\formule{R_i \to R_i + rR_j}$ \\ \hline
\end{tabular}\par}

\begin{stelling}{Gelijkwaardige stelsels}
Een elementaire rijbewerking verandert de oplossingen van een stelsel niet.
\end{stelling}

\begin{bewijs}
Schrijf vergelijking $i$ kort als $L_i = c_i$, met $L_i$ het linkerlid. Neem een oplossing van het oude stelsel. Verwisselen verandert niets aan de vergelijkingen zelf. Uit $L_i = c_i$ volgt $rL_i = rc_i$, en uit $L_i = c_i$ en $L_j = c_j$ volgt
\bewijsregel{L_i + rL_j = c_i + rc_j .}
De oplossing voldoet dus ook aan het nieuwe stelsel. Omgekeerd maak je elke rijbewerking ongedaan met een rijbewerking van dezelfde soort: terug verwisselen, vermenigvuldigen met $\frac1r$, of $rR_j$ weer aftrekken. Met dezelfde redenering is elke oplossing van het nieuwe stelsel dus ook een oplossing van het oude.
\end{bewijs}
\end{theorie}

Stelsels met dezelfde oplossingen heten \textbf{gelijkwaardig}. Het bewijs toont waarom $r$ bij de tweede bewerking verschillend moet zijn van 0: een rij maal 0 kun je niet terugdraaien. Neem $x + y = 2$ en $x - y = 0$, met als enige oplossing $(1, 1)$. Vermenigvuldig je de tweede vergelijking met 0, dan staat er $0 = 0$, en is ook $(2, 0)$ een oplossing: je bent een vergelijking kwijt.

\begin{oefeningen}
\oef{h2-rijbew}
\end{oefeningen}

\subsection{De rijgereduceerde vorm}

Met rijbewerkingen werk je naar een matrix waaruit je de oplossing zo kunt aflezen. In de mooiste vorm staat links de eenheidsmatrix: dan is elke rij een vergelijking van de vorm $x = \ldots$

\begin{theorie}
Een matrix is \textbf{\textcolor{accent}{rijgereduceerd}} als

\vspace{1mm}
\quad 1.\ \ het eerste getal verschillend van 0 in elke rij een 1 is: de \textbf{\textcolor{accent}{spil}} van die rij;

\quad 2.\ \ elke spil rechts staat van de spil in de rij erboven, en de nulrijen onderaan staan;

\quad 3.\ \ in de kolom van een spil verder alleen nullen staan.
\end{theorie}

\figuren{$\um{ccc|c}{\textcolor{accent}{\mathbf{1}} & 0 & 0 & 1 \\ 0 & \textcolor{accent}{\mathbf{1}} & 0 & 2 \\ 0 & 0 & \textcolor{accent}{\mathbf{1}} & 3}$ \qquad\qquad $\um{ccc|c}{\textcolor{accent}{\mathbf{1}} & 0 & -7 & -1 \\ 0 & \textcolor{accent}{\mathbf{1}} & 3 & 2 \\ 0 & 0 & 0 & 0}$ \qquad\qquad $\um{ccc|c}{\textcolor{accent}{\mathbf{1}} & 2 & 0 & 4 \\ 0 & 0 & \textcolor{accent}{\mathbf{1}} & 5 \\ 0 & 0 & 0 & 0}$}

Drie rijgereduceerde matrices, met de spillen in kleur. Links lees je meteen af: $x = 1$, $y = 2$ en $z = 3$. De andere twee hebben een nulrij; wat dat betekent, zie je in sectie 3.

\begin{oefeningen}
\oef{h2-rijgered}
\end{oefeningen}

Elke matrix kun je met rijbewerkingen rijgereduceerd maken, en dat resultaat is uniek: welke bewerkingen je ook kiest, je komt altijd bij dezelfde matrix uit. Dat bewijzen we niet. De werkwijze is vast: kolom per kolom, van links naar rechts.

\voorbeeld{}{Los op met de methode van Gauss-Jordan: $\begin{cases} x + y + z = 6 \\ 2x + 3y + z = 11 \\ x - y + 2z = 5 \end{cases}$\label{vb:h2-gauss}}{%
\uitleg{Dit is het stelsel uit de Verken. We maken eerst in kolom 1 een spil en nullen eronder, dan in kolom 2, dan in kolom 3. Elke stap zetten we naast de matrix.}
\begin{stappenbreed}
\um{ccc|c}{1 & 1 & 1 & 6 \\ 2 & 3 & 1 & 11 \\ 1 & -1 & 2 & 5} & De uitgebreide matrix. Linksboven staat al een 1: de spil van kolom 1. \\[5ex]
\um{ccc|c}{1 & 1 & 1 & 6 \\ 0 & 1 & -1 & -1 \\ 0 & -2 & 1 & -1} \ \begin{array}{l} \\ \rb{R_2 - 2R_1} \\ \rb{R_3 - R_1} \end{array} & Nullen onder de spil van kolom 1. \\[5ex]
\um{ccc|c}{1 & 0 & 2 & 7 \\ 0 & 1 & -1 & -1 \\ 0 & 0 & -1 & -3} \ \begin{array}{l} \rb{R_1 - R_2} \\ \\ \rb{R_3 + 2R_2} \end{array} & In kolom 2 staat de spil al op 1. Nullen erboven en eronder. \\[5ex]
\um{ccc|c}{1 & 0 & 0 & 1 \\ 0 & 1 & 0 & 2 \\ 0 & 0 & 1 & 3} \ \begin{array}{l} \rb{R_1 - 2R_3'} \\ \rb{R_2 + R_3'} \\ \rb{R_3' = -R_3} \end{array} & Maak eerst de spil van kolom 3 gelijk aan 1 ($R_3' = -R_3$), dan nullen erboven. \\[5ex]
x = 1, \quad y = 2, \quad z = 3 & Aflezen. Controle in de derde vergelijking: $1 - 2 + 6 = 5$.
\end{stappenbreed}}

\begin{werkwijze}
Schrijf de uitgebreide matrix. Werk kolom per kolom, van links naar rechts. Zorg voor een getal verschillend van 0 op de plaats van de spil (verwissel zo nodig rijen), en deel de rij door dat getal. Maak met de spilrij alle andere getallen in die kolom 0. Lees op het einde de oplossing af, en controleer in de oorspronkelijke vergelijkingen.
\end{werkwijze}

Staat er op de plaats van een spil een 0, dan verwissel je eerst twee rijen. Je kunt een stap ook overslaan als het getal al 1 is, of breuken vermijden door eerst een andere rij te kiezen: elke weg leidt naar dezelfde rijgereduceerde matrix.

\voorbeeld{}{Los op: $\begin{cases} y + 2z = 5 \\ x + y + z = 4 \\ 2x + y - z = 1 \end{cases}$\label{vb:h2-gauss2}}{%
\uitleg{In de eerste vergelijking ontbreekt $x$: linksboven staat een 0. We verwisselen eerst de eerste twee rijen.}
\begin{stappenbreed}
\um{ccc|c}{1 & 1 & 1 & 4 \\ 0 & 1 & 2 & 5 \\ 0 & -1 & -3 & -7} \ \begin{array}{l} \rb{R_1 \leftrightarrow R_2} \\ \\ \rb{R_3 - 2R_1} \end{array} & Na de verwisseling staat de spil 1 linksboven. \\[5ex]
\um{ccc|c}{1 & 0 & -1 & -1 \\ 0 & 1 & 2 & 5 \\ 0 & 0 & -1 & -2} \ \begin{array}{l} \rb{R_1 - R_2} \\ \\ \rb{R_3 + R_2} \end{array} & Kolom 2. \\[5ex]
\um{ccc|c}{1 & 0 & 0 & 1 \\ 0 & 1 & 0 & 1 \\ 0 & 0 & 1 & 2} & Kolom 3: eerst $R_3 \to -R_3$, dan $R_1 + R_3$ en $R_2 - 2R_3$. \\[5ex]
x = 1, \quad y = 1, \quad z = 2 & Controle: $1 + 4 = 5$, $1 + 1 + 2 = 4$ en $2 + 1 - 2 = 1$.
\end{stappenbreed}}

\nujij{Los op met Gauss-Jordan: $x + 2y = 1$ en $3x + 5y = 4$.}{$\um{cc|c}{1 & 2 & 1 \\ 3 & 5 & 4} \to \um{cc|c}{1 & 2 & 1 \\ 0 & -1 & 1} \to \um{cc|c}{1 & 0 & 3 \\ 0 & 1 & -1}$: $x = 3$, $y = -1$.}

\begin{oefeningen}
\oef{h2-gauss}

\oef{h2-wissel}
\end{oefeningen}

%========================================================================================
\section{Bepaald, onbepaald of strijdig}

\subsection{Drie soorten stelsels}

Bij twee onbekenden is elke vergelijking een rechte. Twee rechten snijden elkaar in één punt, vallen samen, of zijn evenwijdig: één oplossing, oneindig veel, of geen. Dat blijft zo bij meer onbekenden, en de rijgereduceerde matrix toont welk geval het is.

\figuren{\begin{tikzpicture}[x=.42cm,y=.42cm,font=\scriptsize]
  \draw[->,grijs] (-.5,0) -- (4.5,0); \draw[->,grijs] (0,-.5) -- (0,4.5);
  \draw[accent,thick] (0,.5) -- (4,3.5); \draw[accent,thick] (0,4) -- (4,0);
  \fill[accent] (2,2) circle (2pt);
  \node[below,text=black] at (2,-.6) {één oplossing};
\end{tikzpicture}\qquad\qquad
\begin{tikzpicture}[x=.42cm,y=.42cm,font=\scriptsize]
  \draw[->,grijs] (-.5,0) -- (4.5,0); \draw[->,grijs] (0,-.5) -- (0,4.5);
  \draw[accent,line width=2.2pt] (0,1) -- (4,4); \draw[fel!80!black,thick,dashed] (0,1) -- (4,4);
  \node[below,text=black] at (2,-.6) {oneindig veel};
\end{tikzpicture}\qquad\qquad
\begin{tikzpicture}[x=.42cm,y=.42cm,font=\scriptsize]
  \draw[->,grijs] (-.5,0) -- (4.5,0); \draw[->,grijs] (0,-.5) -- (0,4.5);
  \draw[accent,thick] (0,.5) -- (4,3.5); \draw[accent,thick] (0,2) -- (4,5);
  \node[below,text=black] at (2,-.6) {geen};
\end{tikzpicture}}

\begin{theorie}
Een stelsel is \textbf{\textcolor{accent}{strijdig}} als het geen oplossing heeft, \textbf{\textcolor{accent}{bepaald}} als het precies één oplossing heeft, en \textbf{\textcolor{accent}{onbepaald}} als het oneindig veel oplossingen heeft.

\begin{stelling}{Soort stelsel aflezen}
Rijherleid de uitgebreide matrix. Dan is het stelsel

\vspace{1mm}
{\centering\arrayrulecolor{accent}\renewcommand{\arraystretch}{1.5}%
\begin{tabular}{|>{\centering\arraybackslash}m{2.5cm}|>{\centering\arraybackslash}m{9.6cm}|}
\hline
\textbf{\textcolor{accent}{strijdig}} & als er een rij $\formule{[0 \ \cdots \ 0 \,|\, c]}$ is met $c \neq 0$ \\ \hline
\textbf{\textcolor{accent}{bepaald}} & als het niet strijdig is, en elke onbekende een spil heeft \\ \hline
\textbf{\textcolor{accent}{onbepaald}} & als het niet strijdig is, en minstens één onbekende geen spil heeft \\ \hline
\end{tabular}\par}

\vspace{1mm}
Bij een onbepaald stelsel is elke onbekende zonder spil een \textbf{\textcolor{accent}{parameter}}: ze mag elke waarde aannemen.
\end{stelling}

\begin{bewijs}
Rijbewerkingen veranderen de oplossingen niet, dus kijken we naar het rijgereduceerde stelsel. Een rij $[0 \ \cdots \ 0 \,|\, c]$ staat voor $0 = c$. Is $c \neq 0$, dan klopt die vergelijking nooit: geen oplossing. Is er geen zo'n rij, dan zeggen de nulrijen $0 = 0$, en blijft er per spil één vergelijking over. Omdat in de kolom van een spil verder alleen nullen staan, komt in die vergelijking geen andere onbekende met een spil voor:
\bewijsregel{(\text{onbekende met spil}) + (\text{termen met onbekenden zonder spil}) = d .}
Heeft elke onbekende een spil, dan staat er gewoon $x_k = d_k$: precies één oplossing. Zijn er onbekenden zonder spil, geef die dan om het even welke waarden. Elke vergelijking geeft dan de waarde van haar onbekende met spil, en dat is een oplossing. Verschillende waarden geven verschillende oplossingen, dus zijn er oneindig veel.
\end{bewijs}
\end{theorie}

Een nulrij $[0 \ 0 \ 0 \,|\, 0]$ zegt dus niets: die vergelijking was een gevolg van de andere. Pas op: een stelsel met een nulrij is daarom nog niet onbepaald. Dat beslis je met de spillen. Zo geeft het stelsel $x + y = 2$, $x - y = 0$, $2x = 2$ na rijherleiden een nulrij, maar $x$ en $y$ hebben allebei een spil: het is bepaald, met oplossing $(1, 1)$.

\voorbeeld{}{Los op: $\begin{cases} x + 2y - z = 3 \\ 2x + 5y + z = 8 \\ x + 3y + 2z = 5 \end{cases}$\label{vb:h2-onbepaald}}{%
\uitleg{We rijherleiden zoals altijd. Er blijft een nulrij over, en de derde kolom krijgt geen spil: $z$ wordt een parameter.}
\begin{stappenbreed}
\um{ccc|c}{1 & 2 & -1 & 3 \\ 0 & 1 & 3 & 2 \\ 0 & 1 & 3 & 2} \ \begin{array}{l} \\ \rb{R_2 - 2R_1} \\ \rb{R_3 - R_1} \end{array} & De laatste twee rijen zijn gelijk. \\[5ex]
\um{ccc|c}{1 & 0 & -7 & -1 \\ 0 & 1 & 3 & 2 \\ 0 & 0 & 0 & 0} \ \begin{array}{l} \rb{R_1 - 2R_2} \\ \\ \rb{R_3 - R_2} \end{array} & Rijgereduceerd. Geen strijdige rij, en kolom 3 heeft geen spil: onbepaald. \\[5ex]
x - 7z = -1, \quad y + 3z = 2 & De twee overgebleven vergelijkingen. \\[1ex]
(x, y, z) = (-1 + 7t, \ 2 - 3t, \ t), \quad t \in \R & Kies $z = t$ en zonder $x$ en $y$ af.
\end{stappenbreed}
Voor elke $t$ is dat een oplossing: met $t = 0$ is dat $(-1, 2, 0)$, met $t = 1$ is het $(6, -1, 1)$.}

\voorbeeld{}{Los op: $\begin{cases} x + 2y - z = 3 \\ 2x + 5y + z = 8 \\ x + 3y + 2z = 6 \end{cases}$\label{vb:h2-strijdig}}{%
\uitleg{Hetzelfde stelsel, met alleen een ander rechterlid in de derde vergelijking. We doen dezelfde rijbewerkingen.}
\begin{stappenbreed}
\um{ccc|c}{1 & 2 & -1 & 3 \\ 0 & 1 & 3 & 2 \\ 0 & 1 & 3 & 3} \ \begin{array}{l} \\ \rb{R_2 - 2R_1} \\ \rb{R_3 - R_1} \end{array} & Nu verschillen de laatste twee rijen alleen in het rechterlid. \\[5ex]
\um{ccc|c}{1 & 0 & -7 & -1 \\ 0 & 1 & 3 & 2 \\ 0 & 0 & 0 & 1} \ \begin{array}{l} \rb{R_1 - 2R_2} \\ \\ \rb{R_3 - R_2} \end{array} & De laatste rij betekent $0 = 1$: het stelsel is strijdig.
\end{stappenbreed}}

\begin{werkwijze}
Rijherleid de uitgebreide matrix. Is er een rij $[0 \ \cdots \ 0 \,|\, c]$ met $c \neq 0$, dan is het stelsel strijdig. Anders: heeft elke onbekende een spil, dan lees je de enige oplossing af. Zijn er onbekenden zonder spil, kies die dan als parameters en druk de andere onbekenden erin uit.
\end{werkwijze}

\nujij{Hoe ziet de oplossing eruit als de rijgereduceerde matrix $\um{ccc|c}{1 & 2 & 0 & 4 \\ 0 & 0 & 1 & 5 \\ 0 & 0 & 0 & 0}$ is?}{Kolom 2 heeft geen spil: $y = t$. Dan is $x = 4 - 2t$ en $z = 5$: $(4 - 2t, t, 5)$.}

\begin{oefeningen}
\oef{h2-soort}

\oef{h2-soort2}
\end{oefeningen}

\subsection{Stelsels met een parameter}

Staat er in een stelsel een letter die geen onbekende is, dan hangt het soort stelsel soms af van de waarde van die letter. Je rijherleidt zoals gewoonlijk, maar je deelt nooit door een uitdrukking die 0 kan zijn zonder dat geval apart te bekijken.

\voorbeeld{}{Bespreek het stelsel $\begin{cases} x + y + z = 1 \\ x + 2y + 3z = 2 \\ x + 3y + az = 3 \end{cases}$ voor alle reële waarden van $a$.\label{vb:h2-parameter}}{%
\uitleg{We rijherleiden tot de parameter in een spil zou komen. Daar splitsen we: is die uitdrukking 0 of niet?}
\begin{stappenbreed}
\um{ccc|c}{1 & 1 & 1 & 1 \\ 0 & 1 & 2 & 1 \\ 0 & 2 & a - 1 & 2} \ \begin{array}{l} \\ \rb{R_2 - R_1} \\ \rb{R_3 - R_1} \end{array} & Kolom 1. \\[5ex]
\um{ccc|c}{1 & 0 & -1 & 0 \\ 0 & 1 & 2 & 1 \\ 0 & 0 & a - 5 & 0} \ \begin{array}{l} \rb{R_1 - R_2} \\ \\ \rb{R_3 - 2R_2} \end{array} & Kolom 2. De spil van kolom 3 zou $a - 5$ worden. \\[5ex]
a \neq 5: \quad z = 0, \ y = 1, \ x = 0 & Deel $R_3$ door $a - 5 \neq 0$: $z = 0$. Bepaald, met oplossing $(0, 1, 0)$. \\[1.5ex]
a = 5: \quad (x, y, z) = (t, \ 1 - 2t, \ t) & De laatste rij wordt een nulrij: onbepaald, met $z = t$ als parameter.
\end{stappenbreed}
Strijdig is het stelsel nooit: het rechterlid van de laatste rij is altijd 0.}

\begin{werkwijze}
Rijherleid zoals gewoonlijk, maar deel nooit door een uitdrukking met de parameter die 0 kan worden. Zoek de waarden van de parameter waarvoor een spil 0 wordt. Bespreek die waarden elk apart, en alle andere waarden samen.
\end{werkwijze}

\begin{oefeningen}
\oef{h2-parameter}

\oef{h2-parameter2}
\end{oefeningen}

%========================================================================================
\opdrachtenkop

\begin{oefeningen}
\oef{h2-o16}

\needspace{14\baselineskip}
\oef{h2-o15}

\oef{h2-o3}

\oef{h2-o19}

\oef{h2-o2}

\oef{h2-o1}

\needspace{14\baselineskip}
\oef{h2-o18}

\oef{h2-o4}

\needspace{16\baselineskip}
\oef{h2-o17}

\needspace{14\baselineskip}
\oef{h2-o20}

\oef{h2-o5}

\oef{h2-o8}

\needspace{24\baselineskip}
\oef{h2-o22}

\needspace{14\baselineskip}
\oef{h2-o21}

\oef{h2-o9}

\oef{h2-o7}

\oef{h2-o24}

\oef{h2-o23}

\oef{h2-o10}

\needspace{14\baselineskip}
\oef{h2-o6}

\oef{h2-o11}

\oef{h2-o25}

\oef{h2-o12}

\needspace{14\baselineskip}
\oef{h2-o26}

\oef{h2-o13}

\oef{h2-o14}
\end{oefeningen}

%----------------------------------------------------------------------------------------
% KOPPEL: verbanden met andere leerstof
%----------------------------------------------------------------------------------------
\needspace{20\baselineskip}
\opdrachtsectie{Koppel}

\lkr{Tweedegraadsfuncties en rechten in het vlak kennen de leerlingen uit de 2de graad. Vlakken in de ruimte komen in de cursus Ruimtemeetkunde (6de jaar): daar komen stelsels met drie onbekenden terug.}

\begin{oefeningen}
\oef{h2-koppel-1}

\oef{h2-koppel-2}
\end{oefeningen}

%----------------------------------------------------------------------------------------
\opdrachtsectie{Test jezelf}

\begin{oefeningen}
\oef{h2-t1}

\oef{h2-t2}

\oef{h2-t3}

\oef{h2-t4}

\oef{h2-t5}

\oef{h2-t6}

\oef{h2-t7}
\end{oefeningen}

\herhaalschema{%
  \fv{oef:h2-t1}{vb:h2-matrixvorm}{oef:h2-matrixvorm}%
  \fv{oef:h2-t2}{vb:h2-gauss}{oef:h2-gauss}%
  \fv{oef:h2-t3}{vb:h2-onbepaald}{oef:h2-soort}%
  \fv{oef:h2-t4}{vb:h2-strijdig}{oef:h2-soort}%
  \fv{oef:h2-t5}{vb:h2-parameter}{oef:h2-parameter}%
  \fv{oef:h2-t6}{vb:h2-gauss2}{oef:h2-gauss}%
  \fzv{oef:h2-t7}{oef:h2-rijgered}}

\ifzelfstudie
\section*{Maak je eigen samenvatting}
Sluit dit boek en schrijf in je schrift: hoe je een stelsel als matrixvergelijking en als uitgebreide matrix schrijft, de drie elementaire rijbewerkingen en waarom ze de oplossingen niet veranderen, wanneer een matrix rijgereduceerd is, de stappen van de methode van Gauss-Jordan, hoe je herkent of een stelsel bepaald, onbepaald of strijdig is, en hoe je een stelsel met een parameter bespreekt. Vergelijk daarna met de kaders.
\fi
